\scanpage{345}
\mspage{84}
% archive sheet number: 344
\[
\underline{\mathrm{Ch}}_\tau(I_{\tau'};0_{\tau'},1_{\tau'})
\quad(\simeq\underline{\mathrm{Hom}}^{\unsure{\varepsilon|E}}(I_\tau,I_{\tau'})),
\]
mais $\underline{\mathrm{Hom}}_{\underline{\mathrm{Ty}}}(\tau,\tau')$ est la sous-catégorie pleine de cette catégorie formée des \uline{chemins croissants} de $I_{\tau'}$. Le lemme d'homotopie signifie ainsi que cette catégorie $\underline{\mathrm{Hom}}_{\underline{\mathrm{Ty}}}(\tau,\tau')$ est \deleted{\unsure{connexe}} toujours \uline{connexe}. Elle n'est d'ailleurs pas forcément \unsure{$s$-connexe|$0$-connexe}, car elle peut être vide. Notons \unsure{ici}, pour plus \unsure{tard|de clarté}, la

\uline{Proposition 1}~: Soient $\tau,\tau'\in\mathrm{Ty}$\inserted{, \deleted{\unsure{soient}} $r,r'$ leurs \unsure{longueurs utiles|longueur utile}.} Pour que $\mathrm{Hom}_{\underline{\mathrm{Ty}}}(\tau,\tau')$ soit non vide, il faut et il suffit que l'une ou \deleted{\unsure{moins}} \inserted{l'autre} des deux conditions suivantes soit satisfaite~:

a) $r\geq r'+1$

b) $r=r'$, et (si $r\ne0$) \deleted{\unsure{le}} $\tau$ et $\tau'$ ont même direction initiale (de sorte que $\tau_{\mathrm{red}}=\tau'_{\mathrm{red}}$).

(Démonstration $\pm$ immédiate)

Notons que si $\tau,\tau'$ n'ont pas $=$ direction \unsure{initiale|initial}, alors tout morphisme $\tau\mathrel{\unsure{-|\to}}\tau'$ est forcément dégénéré sur le premier intervalle critique de $I_\tau$. \deleted{\unsure{Si nous}} On aura donc $\tau=\Delta_\varepsilon^{\ell_0}\cdot\tau_0$ (\unsure{où} $\ell_0$ est la longueur du premier intervalle critique \unsure{$J$}, \deleted{\unsure{déduit}} $\tau_0$ a direction initiale opposée de celle de $\tau$), et un iso
\[
\underline{\mathrm{Hom}}(\tau,\tau')\simeq\underline{\mathrm{Hom}}(\tau_0,\tau'),
\]
et cette fois $\tau_{\unsure{0}},\tau'$ ont $=$ direction \unsure{initiale|initial}. On supposera donc désormais que \uline{la direction initiale} \uline{de $\tau,\tau'$ est la $=$}. La condition $\underline{\mathrm{Hom}}(\tau,\tau')\ne\varnothing$ équivaut maintenant à $\underline{r\geq r'}$, ce que nous pouvons supposer. Je vais alors décrire un élément canonique
